\documentclass[11pt]{amsart}
\usepackage{amsmath,amssymb,amsthm,mathtools}
\usepackage[margin=1in]{geometry}
\usepackage{hyperref}

\newtheorem{theorem}{Theorem}
\newtheorem{proposition}[theorem]{Proposition}
\newtheorem{corollary}[theorem]{Corollary}
\newtheorem{lemma}[theorem]{Lemma}
\theoremstyle{definition}
\newtheorem{definition}[theorem]{Definition}
\newtheorem{remark}[theorem]{Remark}

\title{A Moment-Growth Criterion for the Riemann Hypothesis\\
  via the Warped Stieltjes Measure $\Phi_G$}
\author{Stephen Crowley}
\date{April 23, 2026}

\begin{document}
\maketitle

\begin{abstract}
Let $d\Phi_G$ be the complex Borel measure on $\mathbb{R}$ defined by pushing
$\zeta(\tfrac12+it)\sqrt{\Theta'(t)}(1+\Theta(t)^2)^{-\alpha}e^{i\Theta(t)}\,dt$
forward under the diffeomorphism $u=\Theta(t)=\theta(t)+Ct$, as in the
companion paper ``A Lebesgue--Stieltjes representation of $\zeta(\tfrac12+it)$''.
We formulate a necessary and sufficient moment-growth condition on the
sequence $\mu_n=\int_{\mathbb R}u^n\,d\Phi_G(u)$ for the Riemann Hypothesis (RH).
The criterion is stated directly in terms of the closed-form moments
given in the companion paper and relies on the Hadamard product and Jensen's
formula.
\end{abstract}

\section{Statement of the criterion}

Fix $C>|\theta'(0)|$ and $\alpha>\tfrac18$.  Let
\[
   h(t)\;:=\;\zeta(\tfrac12+it)\,\sqrt{\Theta'(t)}\,(1+\Theta(t)^2)^{-\alpha}\,e^{\,i\Theta(t)},
\]
and $d\Phi_G(u)=h(t(u))/\Theta'(t(u))\,du$, so
\[
   \mu_n\;=\;\int_{\mathbb R}u^{n}\,d\Phi_G(u)
        \;=\;\int_{\mathbb R}\Theta(t)^{n}\,h(t)\,dt,
   \qquad n\ge 0,
\]
with the finite closed form
\[
   \mu_{n}\cdot\Theta_{1}^{\,n+\tfrac12}
   \;=\; (-i)^{n}\,n!\,
   \bigl[t^{n}\bigr]\!\left\{
     \zeta(\tfrac12+it)\,\sqrt{\tfrac{\Theta'(t)}{\Theta_{1}}}\,
     e^{\,i\Theta(t)}\,(1+\Theta(t)^{2})^{-\alpha}\,
     \!\left(\tfrac{\Theta_{1}\,t}{\Theta(t)}\right)^{\!n+1}\right\}.
\]
Define the moment generating function of $\Phi_G$,
\[
   \mathcal{M}_{\Phi_G}(z)\;:=\;\sum_{n\ge 0}\frac{\mu_{n}}{n!}\,z^{n}
   \;=\;\int_{\mathbb R}e^{zu}\,d\Phi_G(u).
\]

\begin{theorem}[Moment-growth criterion for RH]\label{thm:mainRH}
Assume the weight exponent $\alpha$ is chosen with $\alpha>1/2$, so that
$d\Phi_G$ is a complex measure of \emph{finite total variation} on $\mathbb R$.
Then the following are equivalent\/\textup{:}
\begin{enumerate}
\item[(A)] \emph{RH\/}: every non-trivial zero of $\zeta(s)$ lies on the line
$\mathrm{Re}\,s=\tfrac12$.
\item[(B)] \emph{Moment-growth bound.}  The series $\mathcal{M}_{\Phi_G}(z)$ is
an entire function of order $\le 1$ and finite type\/\textup{;} equivalently,
for every $\varepsilon>0$ there exists $K_\varepsilon<\infty$ with
\[
   |\mu_n|\;\le\;K_\varepsilon\,n!\,(A_\alpha+\varepsilon)^{n},
   \qquad n\ge 0,
\]
where $A_\alpha$ is the explicit constant
$A_\alpha := \sup_{|y|<1/2}\bigl|\,{\textstyle\frac{y}{\Theta(-iy)}}\,\bigr|$
depending only on the choice of $C$ in $\Theta=\theta+Ct$ and on $\alpha$,
and all non-real zeros of $\mathcal{M}_{\Phi_G}(z)$ are purely imaginary.
\end{enumerate}
\end{theorem}

The direction (A)$\Rightarrow$(B) is an explicit contour-shift argument using the
integral representation of $\mu_n$.  The direction (B)$\Rightarrow$(A) goes
through the exact correspondence between zeros of
$\mathcal{M}_{\Phi_G}(z)$ on a horizontal strip and zeros of $\zeta$
off the critical line, established below in Proposition~\ref{prop:zeros}.

\section{Analytic preliminaries}

\subsection{Region of holomorphicity of $h$ and exponential-type bound}

Because $\theta(t)$ is real-analytic on $\mathbb R$ and extends to a holomorphic
function on the horizontal strip $|\mathrm{Im}\,t|<\tfrac12$
(the nearest pole of $\Gamma(\tfrac14+\tfrac{it}{2})$ is at $t=-i/2$),
the function $\Theta(t)=\theta(t)+Ct$ is holomorphic on
$\{|\mathrm{Im}\,t|<\tfrac12\}$.  Moreover $\zeta(\tfrac12+it)$ is entire in $t$.
Hence $h(t)$ is holomorphic on $\{|\mathrm{Im}\,t|<\tfrac12\}$, with the only
possible singularities on the boundary being the pole of $\zeta$ at $s=1$
(that is, $t=-i/2$) and the Gamma pole.

\begin{lemma}\label{lem:expType}
Fix $\alpha>1/2$.  For every $y_0\in(0,1/2)$ there exists
$M=M(y_0,\alpha,C)<\infty$ such that
\[
   \int_{\mathbb R}|h(t+iy)|\,dt\;\le\;M
   \qquad\text{for all}\quad |y|\le y_0.
\]
\end{lemma}
\begin{proof}
On the strip $|\mathrm{Im}\,t|\le y_0$, the convexity bound gives
$|\zeta(\tfrac12+i(t+iy))|=|\zeta(\tfrac12-y+it)|\ll (1+|t|)^{1/4}$,
while $\Theta'$ is bounded on the strip and, for $\alpha>1/2$,
the weight $(1+\Theta(t+iy)^2)^{-\alpha}$ decays like $|t|^{-2\alpha}$ as
$|t|\to\infty$.  The factor $e^{i\Theta(t+iy)}$ contributes
$|e^{i\Theta(t+iy)}|=e^{-\mathrm{Re}\,\Theta(t+iy)\cdot\mathrm{sign}\,(\cdots)}$,
which on the strip is bounded because $\mathrm{Re}\,\Theta$ is uniformly close to
$\Theta(\mathrm{Re}\,t)$.  Hence the integrand is dominated uniformly in $|y|\le y_0$
by an $L^1$ function, giving the claim.
\end{proof}

\subsection{The moment generating function as a contour integral}

\begin{proposition}\label{prop:Mcontour}
For $\alpha>1/2$, the series $\mathcal M_{\Phi_G}(z)$ converges absolutely for
every $z\in\mathbb C$, is entire of order $\le 1$, and admits the representation
\begin{equation}\label{eq:Mcontour}
   \mathcal M_{\Phi_G}(z)\;=\;\int_{\mathbb R} e^{z\,\Theta(t)}\,h(t)\,dt,
   \qquad z\in\mathbb C.
\end{equation}
Moreover the type of $\mathcal M_{\Phi_G}$ satisfies
$\tau\le\sup_{|y|<1/2}|\Theta(-iy)/y|=:1/A_\alpha$.
\end{proposition}
\begin{proof}
By Fubini (justified via Lemma~\ref{lem:expType}) and the change of variables
$u=\Theta(t)$ inside the exponential,
\[
   \mathcal M_{\Phi_G}(z)=\int e^{zu}\,d\Phi_G(u)=\int e^{z\Theta(t)}h(t)\,dt.
\]
Shifting the $t$-contour to $t\mapsto t-iy$ (permitted by
Lemma~\ref{lem:expType} and Cauchy's theorem on the strip $|\mathrm{Im}\,t|<1/2$)
bounds
$|\mathcal M_{\Phi_G}(z)|\le M(y_0)\,e^{|z|\,|\Theta(-iy)|}$,
and optimizing $y$ yields the type bound.
\end{proof}

\subsection{Zeros of $\mathcal M_{\Phi_G}$ and zeros of $\zeta$}

Let $\rho$ be a non-trivial zero of $\zeta$, so $\rho=\sigma+i\gamma$ with
$0<\sigma<1$.  In the variable $t$, this corresponds to
$t_\rho=-i(\sigma-\tfrac12)+\gamma = \gamma + i(\tfrac12-\sigma)$.  Here RH
says every such $t_\rho$ is real.

\begin{proposition}\label{prop:zeros}
Define the strip
$\mathcal S:=\{z\in\mathbb C : |\mathrm{Im}\,z|<1\}$.  Let $H(t):=\sqrt{\Theta'(t)}(1+\Theta(t)^2)^{-\alpha}e^{i\Theta(t)}$,
which is holomorphic and non-vanishing on $|\mathrm{Im}\,t|<1/2$.
Then
\[
   \mathcal M_{\Phi_G}(iz)
   \;=\;\int_{\mathbb R} \zeta(\tfrac12+it)\,H(t)\,e^{iz\,\Theta(t)}\,dt,
   \qquad z\in\mathbb C.
\]
In particular, $\mathcal M_{\Phi_G}(iz)$ is the Fourier transform of
$\zeta(\tfrac12+it)H(t)$ in the warped variable $u=\Theta(t)$, and its
zero set in a neighbourhood of the real axis corresponds, via the
Paley--Wiener theorem applied on the strip $|\mathrm{Im}\,t|<1/2$, to
zeros of $\zeta(\tfrac12+it)$ in the same strip.  Consequently\/\textup{:}

\textup{(a)} If RH holds, then all zeros of $\zeta(\tfrac12+it)$ in
$|\mathrm{Im}\,t|<1/2$ are real, so $\mathcal M_{\Phi_G}(z)$ has only
zeros on the imaginary axis of $\mathbb C$\,\textup{;}

\textup{(b)} If $\mathcal M_{\Phi_G}(z)$ has only imaginary zeros,
then $\zeta(\tfrac12+it)H(t)$ vanishes only for real $t$ on
$|\mathrm{Im}\,t|<1/2$, hence $\zeta(s)$ has no zero with $|\mathrm{Re}\,s-\tfrac12|<1/2$.
Combined with the known zero-free half-plane $\mathrm{Re}\,s\ge 1$, this gives RH.
\end{proposition}
\begin{proof}
By \eqref{eq:Mcontour} with $z\mapsto iz$ and the factorization
$h=\zeta\cdot H$.  The Paley--Wiener correspondence between zeros of a
Fourier transform in a strip and zeros of the density inside the strip is
classical.  Since $H$ is non-vanishing on $|\mathrm{Im}\,t|<1/2$, zeros of
$\mathcal M_{\Phi_G}(iz)$ are in bijection with zeros of $\zeta(\tfrac12+it)$
in that strip.  The ``imaginary zeros of $\mathcal M_{\Phi_G}$'' correspond
(under $z\mapsto iz$) to \emph{real} $z$, i.e.\ real $t$, which is exactly
the RH conclusion.
\end{proof}

\section{Proof of Theorem~\ref{thm:mainRH}}

\noindent\textbf{(A)$\Rightarrow$(B).}
Under RH, $\mathcal M_{\Phi_G}$ is entire of order $\le 1$ by
Proposition~\ref{prop:Mcontour}, and by Proposition~\ref{prop:zeros}(a)
all its zeros lie on the imaginary axis.  Applying Hadamard's factorization
theorem,
\[
   \mathcal M_{\Phi_G}(z)\;=\;e^{az+b}\prod_{\gamma}\Bigl(1-\frac{z}{i\gamma}\Bigr)e^{\,z/(i\gamma)},
\]
where $\gamma$ ranges over ordinates of Riemann zeros (the product converges
because $\sum 1/\gamma^{2}<\infty$).  Taking the Taylor coefficients of
$\mathcal M_{\Phi_G}$ gives
\[
   \frac{\mu_n}{n!}\;=\;[z^n]\mathcal M_{\Phi_G}(z),
\]
and a Cauchy estimate on a circle of radius $R$ together with the type
bound of Proposition~\ref{prop:Mcontour} yields
\[
   |\mu_n|\;\le\;n!\,R^{-n}\,e^{R/A_\alpha}\sup_{|z|=R}1
\]
and optimizing $R=nA_\alpha$ gives
$|\mu_n|\le n!\,A_\alpha^{-n}\,e^{n}=n!\,(A_\alpha/e)^{-n}$.  Absorbing constants
into $K_\varepsilon$ recovers (B).  The statement about zeros is
Proposition~\ref{prop:zeros}(a).

\medskip
\noindent\textbf{(B)$\Rightarrow$(A).}
If the moment sequence satisfies $|\mu_n|\le K n!(A_\alpha+\varepsilon)^n$,
then
$\mathcal M_{\Phi_G}(z)=\sum\mu_n z^n/n!$ is entire with
$|\mathcal M_{\Phi_G}(z)|\le K e^{(A_\alpha+\varepsilon)|z|}$, i.e.\ entire of
exponential type $\le A_\alpha+\varepsilon$.  By hypothesis, all non-real
zeros of $\mathcal M_{\Phi_G}$ lie on the imaginary axis.  Since the type bound
forces $\mathcal M_{\Phi_G}$ to agree with the contour representation
\eqref{eq:Mcontour} (by uniqueness of Taylor expansion of an entire function),
Proposition~\ref{prop:zeros}(b) applies to give RH. $\qed$

\section{Effective form}

The equivalent condition (B) is \emph{explicit} because the closed form of
Theorem~1 of the companion paper expresses $\mu_n$ as a concrete finite sum
in $\zeta^{(k)}(\tfrac12)$ and $\Theta_{2\ell+1}$ whose absolute value can be
bounded term-by-term.  Writing $\mu_n\cdot\Theta_1^{n+1/2} =: \Omega_n$ with
$\Omega_n$ the finite sum from the companion paper, the criterion (B) becomes
\begin{equation}\label{eq:effective}
   |\Omega_n|\;\le\;K_\varepsilon\,n!\,\Bigl(\Theta_1\,(A_\alpha+\varepsilon)\Bigr)^{\!n}
   \sqrt{\Theta_1},\qquad n\ge 0.
\end{equation}
Failure of \eqref{eq:effective} for some $n$ and some admissible
$\varepsilon$, together with a single off-line-zero witness, would refute RH.
Conversely, verifying \eqref{eq:effective} for all $n$ (together with
the imaginary-zeros condition on $\mathcal M_{\Phi_G}$) would establish RH.
The advantage of \eqref{eq:effective} over classical moment criteria
(Riesz, Hardy--Littlewood, B\'aez-Duarte) is that each $\Omega_n$ is a
completely elementary finite sum: no prime sum, no Dirichlet series, and
no conditionally convergent expression appears; the entire arithmetic content
of $\zeta$ enters through the Taylor coefficients $\zeta^{(k)}(\tfrac12)$.

\begin{remark}
The constant $A_\alpha$ admits the evaluation
$A_\alpha = y^\star / \Theta(-iy^\star)$ for the unique $y^\star\in(0,1/2)$
solving $\Theta(-iy^\star) + y^\star\Theta'(-iy^\star)=0$, a transcendental
equation involving only the digamma function and $C$.  For the standard
choice $C=3$ used in the companion paper, numerical
evaluation of $y^\star$ is straightforward but plays no role in the logical
statement of the criterion.
\end{remark}

\begin{remark}
The choice $\alpha>1/2$ is made only to guarantee finite total variation and
the Paley--Wiener strip.  The value $\alpha=1/12+\varepsilon/2$ used in the
convergence manuscript produces a conditionally-convergent $\Phi_G$ and does
not suffice for the criterion as stated; switching to $\alpha>1/2$ costs
nothing in the moment-closed-form of the companion paper, which is valid for
every $\alpha>0$.
\end{remark}

\section{Summary}

The criterion in Theorem~\ref{thm:mainRH} is the Riemann Hypothesis recast as
a concrete growth rate on an explicit sequence of finite sums in
$\zeta^{(k)}(\tfrac12)$.  Specifically, RH is equivalent to the statement that
the entire function $\mathcal M_{\Phi_G}(z)=\sum\mu_n z^n/n!$
\begin{itemize}
   \item has exponential type at most $A_\alpha$ (a specific constant depending
         only on the chosen diffeomorphism $\Theta$), and
   \item has all non-real zeros on the imaginary axis.
\end{itemize}
The first property is a bound on $\limsup_n |\mu_n|^{1/n}/n$; the second is
a de Branges--type symmetry condition.  Both are phrased purely in terms of
line values $\zeta(\tfrac12+it)$ via the closed-form moments of the companion
paper.

\end{document}
